\documentclass[11pt,a4paper]{article}
\usepackage[utf8]{inputenc}
\usepackage[vietnamese]{babel}
\usepackage{amsmath,amssymb}
\usepackage{geometry}
\geometry{a4paper, total={180mm,260mm}, left=15mm, right=15mm, top=18mm, bottom=18mm}
\usepackage{tikz}
\usetikzlibrary{shapes.geometric, arrows.meta, positioning, calc}
\usepackage{tcolorbox}
\tcbuselibrary{skins,breakable}
\usepackage{xcolor}
\usepackage{multicol}
\usepackage{array}
\usepackage{fontawesome}
\usepackage{fancyhdr}

\pagestyle{fancy}
\fancyhf{}
\renewcommand{\headrulewidth}{0.5pt}
\renewcommand{\footrulewidth}{0.5pt}
\fancyhead[L]{\small\bfseries\color{teal}\faBook\ SỔ TAY KIẾN THỨC TOÁN 3}
\fancyhead[R]{\small\bfseries\color{teal}Lớp Toán Thầy Lĩnh -- loptoanthaylinh.vn}
\fancyfoot[C]{\small\color{gray}-- Sổ tay kiến thức: Bài 44. Ôn tập chung học kì 1 --}

\definecolor{primary}{HTML}{005F73}
\definecolor{secondary}{HTML}{0A9396}
\definecolor{accent}{HTML}{EE9B00}
\definecolor{danger}{HTML}{AE2012}
\definecolor{boxbg}{HTML}{F4F6F8}

\begin{document}

% --- BANNER TIÊU ĐỀ ---
\begin{tcolorbox}[
  enhanced,
  colback=white,
  colframe=primary,
  arc=3mm,
  boxrule=1.2pt,
  left=4mm, right=4mm, top=2.5mm, bottom=2.5mm,
  drop shadow=black!10!white
]
  \centering
  {\Large\bfseries\color{primary}\faBookmark\ BÀI 44: ÔN TẬP CHUNG HỌC KÌ 1}\\[1.5mm]
  {\small\itshape\color{black!75}Đại tổng kết toàn bộ chương trình Toán 3 Học kì 1: Số học, Phép tính, Hình học, Đo lường trên 1 trang A4}
\end{tcolorbox}

\vspace{-2mm}


\noindent
\begin{minipage}[t]{0.49\textwidth}\vspace{0pt}
\begin{tcolorbox}[
  enhanced,
  colback=blue!4!white,
  colframe=blue!75!black,
  arc=2mm, boxrule=0.8pt,
  left=2mm, right=2mm, top=1.8mm, bottom=1.8mm,
  title={\footnotesize\bfseries\color{white}\faCalculator\ 1. Mạch Số học \& Phép tính HK1}
]
{\footnotesize
\begin{itemize}\setlength{\itemsep}{1.5pt}\setlength{\leftskip}{-2.5mm}
  \item Thành thạo 4 phép tính trong phạm vi 1 000.
  \item Thuộc lòng bảng cửu chương 2 - 9.
  \item Gấp một số lên nhiều lần ($\times$); Giảm một số đi nhiều lần ($:$).
  \item Tìm một phần mấy của một số: lấy số đó chia cho số phần.
  \item Thứ tự biểu thức: Ngoặc trước $\to$ Nhân chia $\to$ Cộng trừ.
\end{itemize}
}
\end{tcolorbox}
\end{minipage}
\hfill
\begin{minipage}[t]{0.49\textwidth}\vspace{0pt}
\begin{tcolorbox}[
  enhanced,
  colback=teal!4!white,
  colframe=teal!75!black,
  arc=2mm, boxrule=0.8pt,
  left=2mm, right=2mm, top=1.8mm, bottom=1.8mm,
  title={\footnotesize\bfseries\color{white}\faSliders\ 2. Mạch Hình học \& Đo lường HK1}
]
{\footnotesize
\begin{itemize}\setlength{\itemsep}{1.5pt}\setlength{\leftskip}{-2.5mm}
  \item Trung điểm đoạn thẳng chia đôi độ dài.
  \item Đường kính gấp $2$ lần bán kính: $d = 2 \times R$.
  \item Kiểm tra góc vuông bằng thước ê-ke.
  \item Bảng đơn vị: $1\text{ cm} = 10\text{ mm}$, $1\text{ kg} = 1000\text{ g}$, $1\ l = 1000\text{ ml}$.
  \item Nhiệt độ đo bằng độ C ($^\circ\text{C}$).
\end{itemize}
}
\end{tcolorbox}
\end{minipage}

\vspace{-2.5mm}
% --- KHỐI 3: SƠ ĐỒ TƯ DUY TIKZ ---
\noindent
\begin{tcolorbox}[
  enhanced,
  colback=orange!4!white,
  colframe=orange!80!black,
  arc=2mm, boxrule=0.8pt,
  left=2mm, right=2mm, top=1.5mm, bottom=1.5mm,
  title={\footnotesize\bfseries\color{white}\faSitemap\ SƠ ĐỒ TƯ DUY: MA TRẬN 4 TRỤ CỘT KIẾN THỨC TOÁN 3 HK1}
]
\centering
\begin{tikzpicture}[
  box/.style={draw, rounded corners=2pt, align=center, font=\scriptsize\bfseries, inner sep=2.5pt, minimum height=0.75cm},
  arr/.style={-{Stealth[length=2mm]}, thick}
]
  \node[box, fill=purple!15, draw=purple!80!black, text width=2.6cm] (t1) at (0, 0) {BẢNG NHÂN CHIA\\Nhân chia 2 đến 9\\Số tròn chục, trăm};
  \node[box, fill=teal!15, draw=teal!80!black, text width=2.6cm] (t2) at (3.5, 0) {BIỂU THỨC SỐ\\Quy tắc ưu tiên\\Trong ngoặc trước};
  \node[box, fill=blue!12, draw=blue!80!black, text width=2.6cm] (t3) at (7.0, 0) {HÌNH HỌC PHẲNG\\Trung điểm, góc vuông\\Hình tròn, khối hộp};
  \node[box, fill=orange!15, draw=orange!80!black, text width=2.6cm] (t4) at (10.5, 0) {ĐO LƯỜNG\\$mm, g, ml, ^\circ\text{C}$\\Thực hành đo};

  \draw[arr, purple!80!black] (t1) -- (t2);
  \draw[arr, teal!80!black] (t2) -- (t3);
  \draw[arr, blue!80!black] (t3) -- (t4);
\end{tikzpicture}
\end{tcolorbox}

\vspace{-2.5mm}
% --- KHỐI 4: VÍ DỤ & MẸO TRÁNH BẪY ---
\noindent
\begin{minipage}[t]{0.49\textwidth}\vspace{0pt}
\begin{tcolorbox}[
  enhanced,
  colback=green!4!white,
  colframe=green!60!black,
  arc=2mm, boxrule=0.8pt,
  left=2mm, right=2mm, top=1.8mm, bottom=1.8mm,
  title={\footnotesize\bfseries\color{white}\faCheckCircle\ Ví dụ thực hành chuẩn mực}
]
{\footnotesize
\textbf{1. Tính giá trị biểu thức:}\\
$150 + 50 \times 2 = 150 + 100 = \mathbf{250}$.\\
$(150 + 50) \times 2 = 200 \times 2 = \mathbf{400}$.\\
\textbf{2. Bài toán 2 bước tổng hợp HK1:}\\
Có $3$ thùng sách, mỗi thùng có $24$ quyển. Đã chuyển đi $35$ quyển. Còn lại:\\
Tổng số sách: $24 \times 3 = 72$ (quyển).\\
Số sách còn lại: $72 - 35 = \mathbf{37}$ (quyển).
}
\end{tcolorbox}
\end{minipage}
\hfill
\begin{minipage}[t]{0.49\textwidth}\vspace{0pt}
\begin{tcolorbox}[
  enhanced,
  colback=red!4!white,
  colframe=red!70!black,
  arc=2mm, boxrule=0.8pt,
  left=2mm, right=2mm, top=1.8mm, bottom=1.8mm,
  title={\footnotesize\bfseries\color{white}\faExclamationTriangle\ Mẹo tránh bẫy trắc nghiệm}
]
{\footnotesize
\begin{itemize}\setlength{\itemsep}{1.5pt}\setlength{\leftskip}{-2.5mm}
  \item \textbf{Bộ 3 bẫy đề thi HK1 thường gặp:}\\
  1. Nhầm thứ tự tính biểu thức.\\
  2. Bỏ sót số 0 ở thương khi chia.\\
  3. Dừng lại ở bước 1 trong bài toán 2 bước tính.
  \item Luôn dành 3 - 5 phút đọc lại bài và thử lại phép tính trước khi nộp bài thi!
\end{itemize}
}
\end{tcolorbox}
\end{minipage}


\end{document}
